\chapter{Information Theory}

There is a moment, early in learning probability, when you start to feel that a probability distribution
is somehow \emph{carrying} something. A fair coin and a coin that lands heads 99.9\% of the time are both
``random,'' yet they feel profoundly different. Watching the fair coin, you are genuinely in suspense
every flip. Watching the lopsided coin, you can predict ``heads'' over and over and be right almost
always; the outcomes barely tell you anything you did not already expect. One distribution is rich with
uncertainty; the other is nearly empty of it.

Information theory is the branch of probability that takes this feeling and turns it into arithmetic. It
gives us a single number---the \emph{entropy}---that measures how much uncertainty a distribution holds,
how surprised we should expect to be, and (astonishingly) how many bits we would need, on average, to
record a draw from it. From that one idea grows a whole web of quantities: conditional entropy, mutual
information, KL divergence, cross-entropy. These are not exotic curiosities. Cross-entropy is the loss
function that trains essentially every modern classifier. KL divergence is the engine of variational
inference. Mutual information is how we measure dependence when correlation is too crude. And, as we will
see, maximizing likelihood---the workhorse of all of statistics---is secretly the very same thing as
minimizing a cross-entropy.

This chapter was born in 1948, in a single paper by Claude Shannon, and it has the rare quality of being
both deep and \emph{usable on the first day}. Let us build it from the ground up, starting with a single
question: how surprised should you be?

\chapterline

\section{Surprise: measuring how much an outcome tells you}
\label{ch33:surprise}

Imagine a friend texts you the result of an event. Some texts are boring; some are jolting. The text
``the sun rose this morning'' carries essentially no information---you already knew it. The text ``it
snowed in the Sahara today'' is electrifying precisely because it was so unlikely. The lesson is immediate
and intuitive:

\begin{intuitionbox}
Information is the resolution of uncertainty. An outcome is \emph{informative} exactly to the degree that
it was \emph{surprising}---that is, to the degree that it was \emph{improbable}. Certain events tell you
nothing. Rare events tell you a lot.
\end{intuitionbox}

So we want a function $S(p)$, the ``surprise'' of an event of probability $p$, that is large when $p$ is
small and small when $p$ is large. But there are infinitely many decreasing functions. What pins down the
\emph{right} one? Shannon's move was to write down a few properties that any honest measure of surprise
\emph{must} have, and then discover that those properties leave essentially no freedom.

\begin{focusbox}
We demand three things of surprise $S(p)$, where $p \in (0,1]$ is the probability of an event:
\begin{enumerate}
  \item \textbf{Certainty is unsurprising.} $S(1) = 0$. Being told something you already knew for sure
        carries no information.
  \item \textbf{Rarer is more surprising.} $S$ is decreasing in $p$: if $p < q$ then $S(p) > S(q)$.
  \item \textbf{Independent surprises add.} If two \emph{independent} events have probabilities $p$ and
        $q$, learning that both happened should surprise you by the sum of the individual surprises:
        $S(pq) = S(p) + S(q)$.
\end{enumerate}
\end{focusbox}

The third axiom is the heart of it, and it deserves a sentence of defense. If I tell you two
\emph{unrelated} pieces of news---say, that a particular fair coin came up heads (probability $1/2$) and
that a particular fair die came up six (probability $1/6$)---the total ``newness'' should be the
information in the coin \emph{plus} the information in the die. There is no double-counting and no
discount, because the events have nothing to do with each other. Their joint probability is $1/2 \cdot
1/6 = 1/12$, so we are demanding $S(1/12) = S(1/2) + S(1/6)$. Additivity over multiplication of
probabilities: that is a functional equation, and it has a famous solution.

\begin{whybox}
\textbf{Why surprise must be a logarithm.} We want $S(pq) = S(p) + S(q)$ for all $p,q \in (0,1]$. Turn
products into sums and sums stay sums---that is the defining property of the logarithm. Let us prove it is
forced, not just convenient.

Substitute $p = e^{-x}$, $q = e^{-y}$ (so $x,y \ge 0$), and define $g(x) = S(e^{-x})$. The functional
equation becomes
\[
  g(x + y) \;=\; S(e^{-(x+y)}) \;=\; S(e^{-x} e^{-y}) \;=\; S(e^{-x}) + S(e^{-y}) \;=\; g(x) + g(y).
\]
This is \emph{Cauchy's functional equation}, $g(x+y) = g(x) + g(y)$. Together with monotonicity (which we
imposed in axiom 2), its only solutions are the linear ones, $g(x) = c\,x$ for a constant $c$. Unwinding,
$S(p) = g(-\ln p) = -c\,\ln p$. Axiom 2 (decreasing) forces $c > 0$, and axiom 1 ($S(1)=0$) is then
automatic since $\ln 1 = 0$. So
\[
  \boxed{\,S(p) = -c\,\ln p = c\,\ln\frac{1}{p}\,}
\]
is the \emph{unique} measure of surprise, up to the positive scale constant $c$.
\end{whybox}

The constant $c$ is just a choice of \emph{units}, exactly like choosing meters versus feet. Two choices
dominate practice:

\begin{itemize}
  \item Take the logarithm in base $2$: $S(p) = \log_2(1/p)$. The unit is the \textbf{bit}. A fair coin
        ($p = 1/2$) carries $\log_2 2 = 1$ bit of surprise---the canonical ``one yes/no question''
        worth of information. This is the data-science default, and we use it throughout.
  \item Take the natural logarithm: $S(p) = \ln(1/p)$. The unit is the \textbf{nat}. This is common in
        statistics and machine learning, because $\ln$ plays nicely with calculus and with the
        likelihoods we will meet later.
\end{itemize}

To convert, recall $\log_2 x = \ln x / \ln 2$, so one nat $= 1/\ln 2 \approx 1.4427$ bits. We write $\log$
without a base when the base does not matter, and pin it down when it does.

\begin{example}[How surprised is a die?]
\label{ch33:die-surprise}
Roll a fair six-sided die. Each face has probability $1/6$, so seeing any particular face carries
\[
  S(1/6) = \log_2 6 \approx 2.585 \text{ bits}.
\]
Seeing a ``$6$ then a $6$'' on two independent rolls carries $\log_2 36 = 2\log_2 6 \approx 5.17$ bits,
exactly twice as much---additivity in action. Now load the die so that ``six'' has probability $1/2$.
Seeing a six now carries only $\log_2 2 = 1$ bit (it is half-expected), while seeing the rare ``one,''
say with probability $1/10$, carries $\log_2 10 \approx 3.32$ bits. Surprising outcomes shout; expected
ones whisper.
\end{example}

\section{Entropy: the average surprise of a distribution}
\label{ch33:entropy}

Surprise is a property of a single \emph{outcome}. But we usually want to characterize the whole
\emph{distribution}: how uncertain is this random variable, on average, before we see it? The natural
answer is to average the surprise over all outcomes, weighting each by how often it occurs. That average
is the \textbf{entropy}.

\begin{definition}[Shannon entropy]
\label{ch33:def-entropy}
Let $X$ be a discrete random variable with pmf $p(x) = \P(X = x)$ over a finite or countable alphabet
$\mathcal{X}$. The \emph{entropy} of $X$ is the expected surprise,
\[
  H(X) \;=\; \E\!\left[\,S(p(X))\,\right]
        \;=\; \sum_{x \in \mathcal{X}} p(x)\,\log\frac{1}{p(x)}
        \;=\; -\sum_{x \in \mathcal{X}} p(x)\,\log p(x).
\]
We adopt the convention $0\log 0 = 0$ (justified because $p\log p \to 0$ as $p \to 0^+$), so
zero-probability outcomes contribute nothing. With $\log = \log_2$ the units are bits; with $\ln$, nats.
\end{definition}

A crucial point of language: entropy depends only on the \emph{probabilities}, not on the labels of the
outcomes. The faces of a die could be marked with colors, animals, or prime numbers; relabeling does not
change how uncertain the die is. Entropy is a functional of the probability vector $(p_1, \dots, p_n)$,
and we sometimes write $H(p_1, \dots, p_n)$ to emphasize this.

\begin{happeningbox}
\textbf{Entropy is an expectation in disguise.} Notice the formula is literally
$\E[\,\log(1/p(X))\,]$---the expectation of the random variable ``surprise of whatever I drew.'' Everything
you know about expectations applies. This is why entropy behaves so much like a center of mass of
surprise. When one outcome dominates, almost all the weight sits on a small surprise, and the average is
tiny. When the weight is spread out, every outcome is moderately surprising, and the average is large. The
whole subject is, at bottom, expectations of logarithms of probabilities.
\end{happeningbox}

\subsection{The binary entropy function}

The simplest nontrivial case is a single biased coin: $X = 1$ with probability $p$, $X = 0$ with
probability $1-p$. Its entropy is so important it gets its own name and symbol:
\[
  H_b(p) \;=\; -p\log_2 p - (1-p)\log_2(1-p).
\]
Look at the shape of this function (Figure~\ref{ch33:fig-binary-entropy}). It is zero at $p = 0$ and
$p = 1$---a coin that always lands the same way has no uncertainty, and indeed carries zero information
per flip. It rises to a single maximum of exactly $1$ bit at $p = 1/2$, the fair coin, which is the most
uncertain a binary outcome can be. And it is symmetric: a coin biased $90/10$ is exactly as uncertain as
one biased $10/90$, because uncertainty does not care which side we call ``heads.''

\begin{figure}[H]
\centering
\begin{tikzpicture}
\begin{axis}[
    width=11.5cm, height=7cm,
    xlabel={$p$ (probability of heads)},
    ylabel={$H_b(p)$ (bits)},
    domain=0.001:0.999, samples=160,
    xmin=0, xmax=1, ymin=0, ymax=1.08,
    grid=both, grid style={gray!18},
    legend pos=south east, legend cell align=left,
    title={The binary entropy function},
]
\addplot[distblue, very thick] {-x*log2(x) - (1-x)*log2(1-x)};
\addlegendentry{$H_b(p)=-p\log_2 p-(1-p)\log_2(1-p)$}
\addplot[densityred, dashed, thick] coordinates {(0.5,0) (0.5,1)};
\addlegendentry{maximum at $p=\tfrac12$, $H_b=1$ bit}
\end{axis}
\end{tikzpicture}
\caption{Binary entropy. A fair coin ($p=\tfrac12$) is maximally uncertain at $1$ bit; near-deterministic
coins ($p\to 0$ or $p\to 1$) carry almost no uncertainty. The curve is steep near the ends---moving from
$p=0.5$ to $p=0.6$ barely changes entropy, but moving from $p=0.99$ to $p=0.999$ matters a great deal.}
\label{ch33:fig-binary-entropy}
\end{figure}

\begin{example}[Entropy of a fair vs.\ loaded die]
\label{ch33:fair-loaded-die}
A fair six-sided die has $p(x) = 1/6$ for each face, so
\[
  H = \sum_{x=1}^{6} \tfrac16 \log_2 6 = \log_2 6 \approx 2.585 \text{ bits}.
\]
Now load it: $p = (0.5,\,0.1,\,0.1,\,0.1,\,0.1,\,0.1)$. Then
\[
  H = -0.5\log_2 0.5 - 5\,(0.1\log_2 0.1)
    = 0.5 + 5(0.1 \cdot 3.3219)
    = 0.5 + 1.661 = 2.161 \text{ bits}.
\]
The loaded die is \emph{less} uncertain than the fair one, as it should be---you can guess ``six'' and be
right half the time. Concentrating probability always lowers entropy.
\end{example}

\begin{figure}[H]
\centering
\begin{tikzpicture}
\begin{axis}[
    width=12cm, height=6cm,
    ybar, bar width=10pt,
    xlabel={face}, ylabel={probability $p(x)$},
    xmin=0.5, xmax=6.5, ymin=0, ymax=0.58,
    xtick={1,2,3,4,5,6},
    ymajorgrids=true, grid style={gray!18},
    legend pos=north east, legend cell align=left,
    enlarge x limits=0.12,
]
\addplot[draw=distblue, fill=distblue!35] coordinates {(1,0.1667)(2,0.1667)(3,0.1667)(4,0.1667)(5,0.1667)(6,0.1667)};
\addlegendentry{fair: $H=2.585$ bits}
\addplot[draw=densityred, fill=densityred!30] coordinates {(1,0.1)(2,0.1)(3,0.1)(4,0.1)(5,0.1)(6,0.5)};
\addlegendentry{loaded: $H=2.161$ bits}
\end{axis}
\end{tikzpicture}
\caption{Probability mass functions of the fair die (flat, uniform) and the loaded die from
Example~\ref{ch33:fair-loaded-die}, which dumps half its mass on the face ``$6$.'' Concentrating mass on
one outcome makes the die more predictable and lowers its entropy from $2.585$ to $2.161$ bits.}
\label{ch33:fig-die-pmf}
\end{figure}

\subsection{The bounds: why uniform is the most uncertain}

Two facts about entropy are so basic that they anchor all intuition. First, entropy is never negative.
Second, on an alphabet of $n$ symbols it never exceeds $\log n$, with equality exactly for the uniform
distribution.

\begin{theorem}[Bounds on entropy]
\label{ch33:entropy-bounds}
For a discrete random variable $X$ on an alphabet of size $n$,
\[
  0 \;\le\; H(X) \;\le\; \log n,
\]
where the lower bound holds with equality iff $X$ is deterministic (some outcome has probability $1$), and
the upper bound holds with equality iff $X$ is uniform ($p(x) = 1/n$ for all $x$).
\end{theorem}

\begin{whybox}
\textbf{Why uniform maximizes entropy.} The lower bound is easy: each term $p(x)\log(1/p(x)) \ge 0$ since
$0 \le p(x) \le 1$, so the sum is nonnegative, and it vanishes exactly when every term does, i.e.\ when
some $p(x) = 1$.

For the upper bound, the slick proof uses \emph{Jensen's inequality} for the concave function $\log$.
Write the entropy as an expectation and pull the $\log$ outside:
\begin{align*}
  H(X) = \sum_x p(x)\log\frac{1}{p(x)} = \E\!\left[\log\frac{1}{p(X)}\right]
       &\;\le\; \log \E\!\left[\frac{1}{p(X)}\right]
        = \log\!\left(\sum_x p(x)\cdot\frac{1}{p(x)}\right) \\
       &= \log\!\left(\sum_{x:\,p(x)>0} 1\right) \le \log n.
\end{align*}
The single inequality is Jensen's, valid because $\log$ is concave; equality in Jensen requires
$1/p(X)$ to be constant, i.e.\ all the $p(x)$ equal, i.e.\ uniform. The final $\le \log n$ counts the
support. So entropy is maximized by spreading probability as evenly as possible---uncertainty is
greatest when no outcome is favored.
\end{whybox}

\begin{aestheticbox}
This is the ``maximum entropy principle'' in its purest form: among all distributions on $n$ symbols,
the one that assumes the least---commits to nothing beyond ``these $n$ things can happen''---is the
uniform one, and it has the most entropy. Maximizing entropy subject to constraints is a deep modeling
principle: it produces the distribution that is maximally noncommittal given what you actually know. With
a fixed mean it yields the exponential; with a fixed mean and variance it yields the Gaussian. The bell
curve is, in this sense, the most honest distribution with a given spread.
\end{aestheticbox}

\begin{historybox}
\textbf{Shannon, 1948.} Claude Shannon was a research mathematician at Bell Labs when he published
\emph{A Mathematical Theory of Communication} in the \emph{Bell System Technical Journal}. In one stroke
it created the field. The story goes that Shannon was unsure what to call his $-\sum p\log p$ quantity.
John von Neumann reportedly advised him to call it \emph{entropy}---both because the identical formula
appears in statistical mechanics (where Boltzmann and Gibbs had written $S = -k\sum p\log p$ decades
earlier), and, von Neumann supposedly added with a wink, because ``nobody knows what entropy really is,
so in a debate you will always have the advantage.'' The pun is real: thermodynamic entropy and Shannon
entropy are the \emph{same} mathematical object, measuring missing information about microscopic states.
\end{historybox}

\section{Source coding: entropy is the cost of description}
\label{ch33:source-coding}

So far entropy is ``average surprise,'' which is poetic but abstract. Shannon's first great theorem makes
it concrete and a little miraculous: \emph{entropy is the minimum average number of bits needed to
encode draws from a distribution.} It is not a metaphor that entropy is measured in bits. It is the
literal, unbeatable cost of writing the outcomes down.

\begin{intuitionbox}
Suppose you must transmit a long sequence of draws from $X$ over a channel that only speaks $0$s and $1$s,
using as few symbols as possible. The trick is to give \emph{short codewords to frequent outcomes and
long codewords to rare ones}---exactly what Morse code does (``E'' is a single dot; ``Q'' is dash-dash-dot-dash).
The surprise $\log_2(1/p(x))$ turns out to be the ideal length, in bits, for the codeword of an outcome
that occurs with probability $p(x)$. Averaging those ideal lengths gives the entropy.
\end{intuitionbox}

\begin{theorem}[Shannon's source coding theorem, informally]
\label{ch33:source-coding-thm}
To encode i.i.d.\ draws from $X$ with a uniquely decodable binary code, the average codeword length $L$
must satisfy $L \ge H(X)$ (in bits). Conversely, you can always build a code with $L < H(X) + 1$, and by
coding long blocks of draws together you can push the average bits-per-symbol arbitrarily close to
$H(X)$. The entropy is the fundamental limit of lossless compression.
\end{theorem}

\begin{walkbox}
Why can't a code beat the assignment ``length $\approx \log_2(1/p(x))$''? Because of a budget. In any
prefix-free binary code (no codeword is the start of another, so decoding is unambiguous), the codeword
lengths $\ell(x)$ must satisfy \emph{Kraft's inequality} $\sum_x 2^{-\ell(x)} \le 1$---there is only so
much ``room'' in the tree of binary strings. Minimizing the average length $\sum_x p(x)\ell(x)$ subject to
that budget (treat the $\ell$ as continuous and use Lagrange multipliers) gives the optimum $\ell(x) =
\log_2(1/p(x))$, and the minimized average is exactly $\sum_x p(x)\log_2(1/p(x)) = H(X)$. The integer
constraint on real codeword lengths costs at most $1$ extra bit, which is the ``$+1$'' in the theorem, and
it washes out when you code blocks.
\end{walkbox}

\begin{example}[A code that hits the bound]
\label{ch33:huffman-example}
Let $X$ take four values with probabilities $\tfrac12, \tfrac14, \tfrac18, \tfrac18$. The entropy is
\[
  H(X) = \tfrac12\cdot 1 + \tfrac14\cdot 2 + \tfrac18\cdot 3 + \tfrac18\cdot 3
       = 0.5 + 0.5 + 0.375 + 0.375 = 1.75 \text{ bits.}
\]
Now assign codewords by surprise: $a \mapsto \texttt{0}$ (length $1 = \log_2 2$), $b \mapsto \texttt{10}$
(length $2$), $c \mapsto \texttt{110}$, $d \mapsto \texttt{111}$ (length $3$ each). This is prefix-free,
and its average length is $\tfrac12(1) + \tfrac14(2) + \tfrac18(3) + \tfrac18(3) = 1.75$ bits---it
\emph{exactly} matches the entropy, because every probability is a power of $1/2$ so the ideal lengths are
integers. A naive fixed-length code would spend $2$ bits per symbol ($\log_2 4$). Matching code lengths to
surprise saves $0.25$ bits per symbol, forever.
\end{example}

\begin{connectionbox}
The Lagrange-multiplier optimization in the walking question is the same machinery you met in
multivariable calculus: minimize $\sum_x p(x) \ell(x)$ on the constraint surface $\sum_x 2^{-\ell(x)} = 1$. The
gradient condition $\nabla(\text{objective}) = \lambda\,\nabla(\text{constraint})$ delivers $\ell(x) \propto
\log(1/p(x))$ immediately. See the companion calculus volume's chapter on constrained optimization; the
``shadow price'' $\lambda$ here is literally the conversion between probability and bits.
\end{connectionbox}

\section{Joint and conditional entropy}
\label{ch33:joint-conditional}

Real problems involve several random variables at once, and the questions become relational: how much
total uncertainty is in the pair $(X,Y)$? Once I learn $Y$, how much uncertainty about $X$ remains? These
are answered by joint and conditional entropy, and they obey a beautifully simple bookkeeping.

\begin{definition}[Joint and conditional entropy]
\label{ch33:def-joint-cond}
For discrete $X,Y$ with joint pmf $p(x,y)$, the \emph{joint entropy} is
\[
  H(X,Y) = -\sum_{x,y} p(x,y)\log p(x,y),
\]
simply the entropy of the pair viewed as one combined variable. The \emph{conditional entropy} of $X$
given $Y$ is the average, over $y$, of the remaining uncertainty about $X$:
\[
  H(X \given Y) = \sum_y p(y)\,H(X \given Y = y)
              = -\sum_{x,y} p(x,y)\log p(x \given y).
\]
\end{definition}

The conditional entropy $H(X \given Y)$ answers exactly the question ``after I am told $Y$, how surprised
will I still be by $X$, on average?'' Note the averaging: for a \emph{particular} value $Y=y$ the residual
uncertainty $H(X \given Y=y)$ could go up or down, but \emph{averaged} over $Y$, conditioning can only
help---more on that shortly.

\begin{theorem}[Chain rule for entropy]
\label{ch33:chain-rule}
\[
  H(X,Y) = H(X) + H(Y \given X) = H(Y) + H(X \given Y).
\]
\end{theorem}

\begin{whybox}
\textbf{Why the chain rule is just $\log$ of the multiplication rule.} The whole identity falls out of
$p(x,y) = p(x)\,p(y \given x)$, the multiplication rule for probabilities, after taking $-\log$ and
averaging. Watch:
\[
  H(X,Y) = -\sum_{x,y} p(x,y)\log p(x,y)
         = -\sum_{x,y} p(x,y)\big[\log p(x) + \log p(y\given x)\big].
\]
Split the sum. The first piece is $-\sum_{x,y} p(x,y)\log p(x) = -\sum_x p(x)\log p(x) = H(X)$ (we summed
out $y$). The second piece is exactly $-\sum_{x,y} p(x,y)\log p(y\given x) = H(Y\given X)$. So $H(X,Y) =
H(X) + H(Y\given X)$. The symmetric version comes from factoring $p(x,y) = p(y)p(x\given y)$ instead.
\end{whybox}

\begin{happeningbox}
The chain rule says uncertainty is additive in the right currency. The total surprise of the pair equals
the surprise of $X$ \emph{plus} whatever surprise about $Y$ is left after $X$ is known. It mirrors exactly
how you would describe the pair to someone: first tell them $X$ (costs $H(X)$ bits), then tell them $Y$
given that they now know $X$ (costs $H(Y\given X)$ bits). Information bookkeeping is just careful
double-entry accounting of surprise.
\end{happeningbox}

\begin{theorem}[Conditioning reduces entropy]
\label{ch33:cond-reduces}
$H(X \given Y) \le H(X)$, with equality iff $X \indep Y$.
\end{theorem}

This is the formal version of ``information never hurts.'' On average, learning $Y$ cannot increase your
uncertainty about $X$. We will prove it in a moment as a one-line corollary of the nonnegativity of
mutual information. First we need the tool that makes everything click: KL divergence.

\section{Relative entropy (KL divergence): the cost of being wrong}
\label{ch33:kl}

Here is a question that turns out to organize the entire subject. Suppose the data really come from a
distribution $p$, but you, mistakenly, believe they come from $q$. You design your code, your bets, your
model around $q$. How much do you pay for the discrepancy? The answer is the \textbf{Kullback--Leibler
divergence}, also called \emph{relative entropy}.

\begin{definition}[KL divergence]
\label{ch33:def-kl}
For two pmfs $p, q$ on the same alphabet $\mathcal{X}$, the KL divergence \emph{from $q$ to $p$} (the
extra cost of using $q$ when the truth is $p$) is
\[
  \KL(p \,\|\, q) \;=\; \sum_{x} p(x)\,\log\frac{p(x)}{q(x)}
                 \;=\; \E_{X \sim p}\!\left[\log\frac{p(X)}{q(X)}\right].
\]
We require $q(x) = 0 \implies p(x) = 0$ (else the divergence is $+\infty$): you cannot afford a model that
assigns zero probability to something that actually happens.
\end{definition}

\begin{intuitionbox}
Read it as a difference of surprises. Under your wrong model $q$, an outcome $x$ surprises you by
$\log(1/q(x))$ bits. Under the truth $p$, it ``should'' surprise you by $\log(1/p(x))$. The per-outcome
overpayment is $\log(1/q(x)) - \log(1/p(x)) = \log(p(x)/q(x))$, and KL is the average overpayment, weighted
by how often each outcome \emph{really} occurs (weighted by $p$, not $q$, because reality runs on $p$).
\end{intuitionbox}

The single most important fact about KL divergence is that it is never negative, and is zero only when
your model is exactly right. This is \emph{Gibbs' inequality}, and it is the load-bearing wall of the
whole theory.

\begin{theorem}[Gibbs' inequality / nonnegativity of KL]
\label{ch33:gibbs}
$\KL(p\,\|\,q) \ge 0$, with equality iff $p = q$ everywhere.
\end{theorem}

\begin{whybox}
\textbf{Why KL is nonnegative.} Again Jensen, applied to the concave $\log$. Write the \emph{negative}
divergence and pull the log out:
\[
  -\KL(p\,\|\,q) = \sum_x p(x)\log\frac{q(x)}{p(x)}
               = \E_{X\sim p}\!\left[\log\frac{q(X)}{p(X)}\right]
               \;\le\; \log\,\E_{X\sim p}\!\left[\frac{q(X)}{p(X)}\right]
               = \log\sum_x q(x) = \log 1 = 0.
\]
So $-\KL \le 0$, i.e.\ $\KL \ge 0$. Equality in Jensen for the strictly concave $\log$ requires the ratio
$q(X)/p(X)$ to be constant; since both are distributions summing to $1$, that constant must be $1$, forcing
$p = q$. A perfectly correct model pays nothing; any error costs strictly positive divergence.
\end{whybox}

\begin{trapbox}
\textbf{KL is not a distance.} It is tempting to read $\KL(p\,\|\,q)$ as ``how far apart $p$ and $q$ are,''
but it is \emph{not symmetric}: $\KL(p\,\|\,q) \ne \KL(q\,\|\,p)$ in general, and it violates the triangle
inequality. The asymmetry is meaningful, not a defect. $\KL(p\,\|\,q)$ weights errors by the \emph{true}
$p$, so it punishes you for assigning low $q$ to outcomes that truly occur, but barely notices if you waste
$q$-mass on outcomes that never happen. Swapping the arguments swaps which mistake is expensive. In
variational inference, choosing to minimize $\KL(q\,\|\,p)$ versus $\KL(p\,\|\,q)$ gives qualitatively
different approximations (``mode-seeking'' versus ``mean-seeking''). Always mind the order.
\end{trapbox}

\begin{example}[Two coins]
\label{ch33:kl-coins}
Let the truth be a fair coin $p = (0.5, 0.5)$ and your model a biased guess $q = (0.9, 0.1)$. Then
\[
  \KL(p\,\|\,q) = 0.5\log_2\frac{0.5}{0.9} + 0.5\log_2\frac{0.5}{0.1}
               = 0.5(-0.848) + 0.5(2.322) = 0.737 \text{ bits.}
\]
You overpay about $0.74$ bits per flip for believing the coin is biased when it is fair. Reverse the
roles---truth $q = (0.9,0.1)$, model $p = (0.5,0.5)$:
\[
  \KL(q\,\|\,p) = 0.9\log_2\frac{0.9}{0.5} + 0.1\log_2\frac{0.1}{0.5}
               = 0.9(0.848) + 0.1(-2.322) = 0.531 \text{ bits.}
\]
Different number---concrete proof of asymmetry.
\end{example}

\section{Cross-entropy: the loss that trains your models}
\label{ch33:cross-entropy}

KL divergence has a twin that is even more famous in practice, because it is the literal loss function
minimized by nearly every classifier you have ever heard of. It is the \textbf{cross-entropy}.

\begin{definition}[Cross-entropy]
\label{ch33:def-cross-entropy}
The cross-entropy of the model $q$ relative to the truth $p$ is the average surprise you experience when
reality is $p$ but you reason with $q$:
\[
  H(p, q) \;=\; -\sum_x p(x)\log q(x) \;=\; \E_{X\sim p}\!\left[\log\frac{1}{q(X)}\right].
\]
\end{definition}

The three quantities entropy, cross-entropy, and KL are bound by one short, illuminating identity. Split
$\log(p/q) = \log(1/q) - \log(1/p)$ inside the KL sum:
\[
  \KL(p\,\|\,q) = \underbrace{-\sum_x p(x)\log q(x)}_{H(p,q)} \;-\; \underbrace{\Big(-\sum_x p(x)\log p(x)\Big)}_{H(p)},
\]
so
\[
  \boxed{\,H(p, q) = H(p) + \KL(p \,\|\, q)\,}.
\]

\begin{aestheticbox}
Read this identity aloud and you understand all of supervised learning. The cost of coding reality with
your model equals the \emph{irreducible} cost (the true entropy $H(p)$, which no model can avoid) plus the
\emph{avoidable} cost (the divergence $\KL(p\,\|\,q)$, which your error inflicts). Since $H(p)$ does not
depend on your model $q$, \emph{minimizing cross-entropy over $q$ is exactly the same as minimizing KL
divergence over $q$.} Training a model to minimize cross-entropy is training it to make $q$ match the true
$p$ as closely as KL can measure. And because $\KL \ge 0$, cross-entropy is bounded below by $H(p)$, with
the floor reached only by a perfect model.
\end{aestheticbox}

\begin{applicationbox}
\textbf{Why cross-entropy is \emph{the} classification loss.} In multiclass classification, each training
example $i$ has a true label, which we write as a one-hot vector $p^{(i)}$ (all mass on the correct class
$y_i$). The network outputs predicted class probabilities $q^{(i)}$ (the softmax layer). The cross-entropy
loss on example $i$ is
\[
  H(p^{(i)}, q^{(i)}) = -\sum_c p^{(i)}_c \log q^{(i)}_c = -\log q^{(i)}_{y_i},
\]
because the one-hot $p^{(i)}$ kills every term except the true class. So the loss is simply the surprise
the model assigns to the correct answer: $-\log(\text{predicted probability of the right class})$. Confident
and correct $\Rightarrow$ near $0$ loss; confident and wrong $\Rightarrow$ enormous loss. Averaged over the
dataset, minimizing this is minimizing the empirical cross-entropy between the data distribution and the
model---and, by the identity above and the next section, it is \emph{exactly maximum likelihood}. Every
time you call \texttt{CrossEntropyLoss} or \texttt{categorical\_crossentropy}, this is the math running
underneath.
\end{applicationbox}

\begin{figure}[H]
\centering
\begin{tikzpicture}
\begin{axis}[
    width=11.5cm, height=6.5cm,
    xlabel={$q$ = model's predicted probability of the true class},
    ylabel={loss $= -\log_2 q$ (bits)},
    domain=0.02:1, samples=160,
    xmin=0, xmax=1, ymin=0, ymax=6,
    grid=both, grid style={gray!18},
    title={Cross-entropy loss for one correctly labeled example},
]
\addplot[densityred, very thick] {-log2(x)};
\addplot[gray, dashed] coordinates {(1,0) (1,6)};
\node[anchor=west] at (axis cs:0.45,4.3) {\small confident \& wrong: huge penalty};
\node[anchor=east] at (axis cs:0.97,0.6) {\small confident \& right: $\approx 0$};
\end{axis}
\end{tikzpicture}
\caption{The per-example cross-entropy loss $-\log q$ as a function of the probability $q$ the model
assigned to the correct class. The penalty is gentle near $q=1$ but explodes as $q \to 0$: a model that is
confidently wrong is punished savagely, which is exactly the gradient signal that drives learning.}
\label{ch33:fig-crossentropy-loss}
\end{figure}

\section{Maximum likelihood is minimum cross-entropy}
\label{ch33:mle-connection}

We promised that the workhorse of statistics is secretly an information-theoretic quantity. Here is the
payoff. Suppose you observe i.i.d.\ data $x_1, \dots, x_N$ and fit a model $q_\theta$ by maximum
likelihood. The log-likelihood, scaled by $1/N$, is
\[
  \frac{1}{N}\sum_{i=1}^N \log q_\theta(x_i)
  \;=\; \sum_{x} \hat p(x)\,\log q_\theta(x),
\]
where $\hat p(x) = \frac1N \#\{i : x_i = x\}$ is the \emph{empirical distribution} of the data. The
right-hand side is precisely $-H(\hat p, q_\theta)$, the negative cross-entropy between the data
distribution and the model. Therefore
\[
  \underbrace{\max_\theta \frac1N \sum_i \log q_\theta(x_i)}_{\text{maximum likelihood}}
  \;=\;
  \underbrace{\min_\theta\; H(\hat p, q_\theta)}_{\text{minimum cross-entropy}}
  \;=\;
  \underbrace{\min_\theta\; \KL(\hat p \,\|\, q_\theta)}_{\text{best fit in KL}}
  \;+\; \text{const}.
\]

\begin{happeningbox}
\textbf{What maximizing likelihood is really doing.} You are choosing the model $q_\theta$ that
is \emph{closest in KL divergence} to the empirical distribution of your data. As the sample grows,
$\hat p \to p$ (the true distribution) by the law of large numbers, so maximum likelihood drives your
model toward the $q_\theta$ that minimizes $\KL(p\,\|\,q_\theta)$---the best possible approximation to
reality within your model family. This is exactly why MLE is consistent: if the truth is \emph{in} your
family, the KL-closest model is the truth itself, and you find it. Likelihood, cross-entropy, and KL are
three views of one optimization.
\end{happeningbox}

\begin{connectionbox}
This recasts the estimation chapter (Maximum Likelihood) in information-theoretic clothing. There you saw
MLE as ``the parameter that makes the data most probable.'' Here you see it as ``the parameter whose model
distribution is informationally closest to the data.'' Both are correct; the second explains \emph{why}
MLE generalizes, why cross-entropy is the right loss, and why a misspecified model still converges to its
KL-projection onto the truth. The Bayesian-inference chapter's variational methods minimize a KL
divergence directly, making the link explicit.
\end{connectionbox}

\section{Mutual information: how much $Y$ tells you about $X$}
\label{ch33:mutual-info}

We come to the crown jewel. Correlation measures linear co-movement and misses everything nonlinear.
\emph{Mutual information} measures \emph{all} statistical dependence---it is zero exactly when two variables
are independent, no matter how twisted their relationship.

\begin{definition}[Mutual information]
\label{ch33:def-mi}
The mutual information between $X$ and $Y$ is the KL divergence between their joint distribution and the
product of their marginals:
\[
  I(X; Y) \;=\; \KL\big(p(x,y) \,\big\|\, p(x)p(y)\big)
          \;=\; \sum_{x,y} p(x,y)\,\log\frac{p(x,y)}{p(x)\,p(y)}.
\]
\end{definition}

The definition is gorgeous once you see it. Independence \emph{means} $p(x,y) = p(x)p(y)$. So mutual
information literally measures how far the joint distribution is from the ``independent'' distribution
it would be if $X$ and $Y$ ignored each other. By Gibbs' inequality, $I(X;Y) \ge 0$, with equality iff
$p(x,y) = p(x)p(y)$, i.e.\ iff $X \indep Y$. \emph{Mutual information is zero exactly when the variables
are independent}---a clean ``yes/no'' test for dependence of any kind.

\begin{theorem}[Mutual information in terms of entropy]
\label{ch33:mi-entropy}
\[
  I(X;Y) = H(X) - H(X\given Y) = H(Y) - H(Y\given X) = H(X) + H(Y) - H(X,Y).
\]
\end{theorem}

\begin{whybox}
\textbf{Why these forms agree.} Start from the definition and split the log using $p(x,y) = p(x\given
y)p(y)$:
\[
  I(X;Y) = \sum_{x,y} p(x,y)\log\frac{p(x\given y)}{p(x)}
         = \sum_{x,y}p(x,y)\log p(x\given y) - \sum_{x,y}p(x,y)\log p(x).
\]
The first sum is $-H(X\given Y)$; the second is $-H(X)$ (sum out $y$). Hence $I(X;Y) = H(X) -
H(X\given Y)$. By symmetry of the definition in $x$ and $y$, also $I(X;Y) = H(Y) - H(Y\given X)$. Finally,
substitute the chain rule $H(X\given Y) = H(X,Y) - H(Y)$ into the first form to get $I(X;Y) = H(X) + H(Y)
- H(X,Y)$.
\end{whybox}

The middle form $I(X;Y) = H(X) - H(X\given Y)$ is the one to memorize as a sentence: \emph{mutual
information is the reduction in your uncertainty about $X$ once you learn $Y$.} It is the average number of
bits that knowing $Y$ saves you when you have to describe $X$. And since $I(X;Y) \ge 0$, this rearranges
instantly to $H(X\given Y) \le H(X)$---the promised proof that \emph{conditioning reduces entropy}.
Information never hurts, on average.

\begin{figure}[H]
\centering
\begin{tikzpicture}[scale=1.0]
  \begin{scope}
    \fill[sourceblue!22] (0,0) circle (1.85cm);
    \fill[targetgreen!22] (2.0,0) circle (1.85cm);
    \begin{scope}
      \clip (0,0) circle (1.85cm);
      \fill[structurepurple!35] (2.0,0) circle (1.85cm);
    \end{scope}
  \end{scope}
  \draw[sourceblue, thick] (0,0) circle (1.85cm);
  \draw[targetgreen, thick] (2.0,0) circle (1.85cm);
  \node[sourceblue!60!black] at (-1.45,1.75) {$H(X)$};
  \node[targetgreen!60!black] at (3.45,1.75) {$H(Y)$};
  \node[sourceblue!50!black] at (-0.55,0) {\small $H(X\given Y)$};
  \node[targetgreen!50!black] at (2.55,0) {\small $H(Y\given X)$};
  \node[structurepurple!70!black, align=center] at (1.0,0) {\small $I(X;Y)$};
  \node[align=center] at (1.0,-2.55) {\small the union of both circles is $H(X,Y)$};
\end{tikzpicture}
\caption{The information ``Venn diagram.'' Each circle is the entropy of one variable. Their overlap is
the mutual information $I(X;Y)$, the shared uncertainty resolved by either variable. The crescents are the
conditional entropies---what remains uncertain about one variable after the other is known. The total area
of the union is the joint entropy $H(X,Y) = H(X) + H(Y) - I(X;Y)$.}
\label{ch33:fig-venn}
\end{figure}

\begin{example}[A noisy channel]
\label{ch33:noisy-channel}
Send a fair bit $X \in \{0,1\}$ through a channel that flips it with probability $f = 0.1$ (the ``binary
symmetric channel,'' sketched in Figure~\ref{ch33:fig-bsc}). The receiver sees $Y$. By symmetry $Y$ is
also a fair bit, so $H(Y) = 1$. Given the
input, the only uncertainty in the output is the flip, so $H(Y\given X) = H_b(0.1) = -0.1\log_2 0.1 -
0.9\log_2 0.9 \approx 0.469$ bits. Therefore
\[
  I(X;Y) = H(Y) - H(Y\given X) \approx 1 - 0.469 = 0.531 \text{ bits.}
\]
Each use of this noisy channel conveys about half a bit of genuine information about the input. A noiseless
channel ($f=0$) would convey the full bit; a useless channel that flips with probability $0.5$ would
convey $1 - H_b(0.5) = 0$ bits---pure noise tells you nothing. This number, maximized over input
distributions, is the channel's \emph{capacity}, the ceiling on reliable communication rate that Shannon's
second great theorem says you can actually approach.
\end{example}

\begin{figure}[H]
\centering
\begin{tikzpicture}[>=Latex, node distance=3.4cm, every node/.style={font=\small}]
  \node[circle, draw=sourceblue, thick, minimum size=0.9cm] (x0) {$0$};
  \node[circle, draw=sourceblue, thick, minimum size=0.9cm, below=1.8cm of x0] (x1) {$1$};
  \node[circle, draw=targetgreen, thick, minimum size=0.9cm, right=3.4cm of x0] (y0) {$0$};
  \node[circle, draw=targetgreen, thick, minimum size=0.9cm, right=3.4cm of x1] (y1) {$1$};
  \draw[->, sourceblue, thick] (x0) -- node[above]{$1-f$} (y0);
  \draw[->, sourceblue, thick] (x1) -- node[below]{$1-f$} (y1);
  \draw[->, trapred, thick] (x0) -- node[pos=0.25, right]{$f$} (y1);
  \draw[->, trapred, thick] (x1) -- node[pos=0.25, right]{$f$} (y0);
  \node[above=0.2cm of x0] {input $X$};
  \node[above=0.2cm of y0] {output $Y$};
\end{tikzpicture}
\caption{The binary symmetric channel. Each transmitted bit arrives unchanged with probability $1-f$ and
flipped with probability $f$. With $f=0.1$ and a fair input, each use carries
$I(X;Y)=1-H_b(0.1)\approx 0.531$ bits; at $f=0.5$ the two red and blue arrows are equally likely, the
output is independent of the input, and the channel conveys nothing.}
\label{ch33:fig-bsc}
\end{figure}

\begin{figure}[H]
\centering
\begin{tikzpicture}
\begin{axis}[
    width=11.5cm, height=6cm,
    xlabel={flip probability $f$},
    ylabel={$I(X;Y)=1-H_b(f)$ (bits)},
    domain=0.0005:0.9995, samples=160,
    xmin=0, xmax=1, ymin=0, ymax=1.05,
    grid=both, grid style={gray!18},
    title={Information per use of the binary symmetric channel},
]
\addplot[structurepurple, very thick] {1 - (-x*log2(x) - (1-x)*log2(1-x))};
\addplot[gray, dashed] coordinates {(0.5,0) (0.5,1.05)};
\node[anchor=south] at (axis cs:0.5,0.02) {\small useless at $f=\tfrac12$};
\end{axis}
\end{tikzpicture}
\caption{Mutual information of the binary symmetric channel as the flip probability $f$ varies. It is a
full bit at the noiseless ends $f=0$ and $f=1$ (a guaranteed flip is just as informative as no flip, since
the receiver can invert it), and drops to $0$ exactly at $f=\tfrac12$, where output and input are
independent.}
\label{ch33:fig-bsc-capacity}
\end{figure}

\begin{applicationbox}
\textbf{Mutual information in data science.} Because $I(X;Y) = 0$ iff $X \indep Y$ and it captures
nonlinear dependence, mutual information is a go-to tool for \emph{feature selection}: rank candidate
features by how much information they carry about the target, keeping high-$I$ features even when their
\emph{correlation} with the target is zero (think of $Y = X^2$ with symmetric $X$: correlation $0$, mutual
information large). Decision trees split on the feature with the highest \emph{information gain}, which is
exactly a conditional mutual information. And in modern representation learning, objectives like InfoNCE
\emph{maximize} a lower bound on the mutual information between an input and its learned embedding, forcing
the representation to retain what matters.
\end{applicationbox}

\begin{paradoxbox}
\textbf{Can knowing more ever make you more uncertain?} We proved $H(X\given Y) \le H(X)$: conditioning
reduces entropy. Yet everyday experience says ``the more I learned, the more confused I got!'' Both are
true, and the resolution is the word \emph{average}. The theorem bounds the \emph{expected} residual
uncertainty $H(X\given Y) = \sum_y p(y)H(X\given Y=y)$. For a \emph{particular} observed value $Y=y$, the
specific conditional entropy $H(X\given Y=y)$ can absolutely \emph{exceed} $H(X)$. Example: you are fairly
sure a die is fair (low uncertainty about which loading it has), then you observe three sixes in a row
(an unusual $y$); now you genuinely do not know whether it is fair or loaded---your uncertainty has gone
\emph{up}. The theorem guarantees only that across all possible observations, weighted by how likely they
are, the surprising-and-confusing $y$'s are outnumbered by the clarifying ones. Information never hurts
\emph{on average}; any single message certainly can.
\end{paradoxbox}

\section{Continuous variables: differential entropy in one breath}
\label{ch33:differential}

Everything above was discrete. For a continuous random variable with density $f$, the natural analogue is
the \emph{differential entropy}
\[
  h(X) = -\int_{-\infty}^{\infty} f(x)\log f(x)\,\dd x = \E[-\log f(X)],
\]
replacing the sum by an integral. It is genuinely useful---among all densities with a fixed variance
$\sigma^2$, the Gaussian uniquely maximizes differential entropy, at $h = \tfrac12\log(2\pi e\,\sigma^2)$,
which is the maximum-entropy justification for the bell curve. But beware: differential entropy is not a
clean limit of discrete entropy (finely discretizing a continuous variable gives entropy that blows up
like $\log(1/\Delta)$), and it can be negative or change under a change of variables. The quantities that
\emph{do} behave well in the continuous world are the \emph{relative} ones---KL divergence and mutual
information---because the troublesome infinities cancel between the two distributions being compared. That
is one more reason KL and mutual information, not entropy itself, are the durable currency of modern
machine learning.

\begin{connectionbox}
The differential entropy is an honest Riemann/Lebesgue integral $\E[-\log f(X)]$; see the companion
calculus volume for integration of densities, and the measure-theory chapters here for the careful
statement that KL divergence is well defined whenever $p$ is absolutely continuous with respect to $q$,
with the integrand $\log(\dd p/\dd q)$ being the log of the Radon--Nikodym derivative.
\end{connectionbox}

\section{Bringing it together}
\label{ch33:summary}

One idea generated this entire chapter: \emph{surprise is $\log(1/p)$, and information theory is the study
of its averages.}

\begin{itemize}
  \item \textbf{Entropy} $H(X) = \E[\log(1/p(X))]$ is average surprise---uncertainty, and the
        unbeatable cost in bits of describing $X$. It is maximized by the uniform distribution.
  \item \textbf{Conditional entropy} $H(X\given Y)$ is leftover uncertainty after seeing $Y$; the
        \textbf{chain rule} $H(X,Y) = H(X) + H(Y\given X)$ is $\log$ of the multiplication rule.
  \item \textbf{KL divergence} $\KL(p\|q) = \E_p[\log(p/q)] \ge 0$ is the cost of modeling truth $p$ with
        wrong model $q$; it is zero iff $p=q$ (Gibbs).
  \item \textbf{Cross-entropy} $H(p,q) = H(p) + \KL(p\|q)$ is the ML loss; minimizing it is minimizing KL
        is maximizing likelihood.
  \item \textbf{Mutual information} $I(X;Y) = \KL(p_{XY}\|p_Xp_Y) = H(X) - H(X\given Y) \ge 0$ measures
        all dependence and is zero iff $X\indep Y$.
\end{itemize}

Every modern loss function, dependence measure, and compression scheme you will meet is one of these five
quantities wearing a costume. Learn to recognize $\log(1/p)$ underneath, and information theory stops being
a list of formulas and becomes a single, simple story about surprise.

\section{Exercises}
\label{ch33:exercises}

\begin{exercise}[Warm-up entropies]
\label{ch33:ex-warmup}
Compute the entropy in bits of: (a) a fair coin; (b) a fair $8$-sided die; (c) a coin with $p=0.25$;
(d) the distribution $(\tfrac12, \tfrac14, \tfrac18, \tfrac1{16}, \tfrac1{16})$. For (d), compare your
answer to $\log_2 5$ and explain why it is smaller.
\end{exercise}

\begin{exercise}[Surprise adds]
\label{ch33:ex-surprise-adds}
You draw a card from a shuffled standard deck. Express in bits the surprise of the events ``it is a
spade,'' ``it is an ace,'' and ``it is the ace of spades.'' Verify numerically that the third surprise
equals the sum of the first two, and explain which property of independence makes this work.
\end{exercise}

\begin{exercise}[The maximum is uniform]
\label{ch33:ex-max-uniform}
Using Lagrange multipliers, maximize $H(p_1,\dots,p_n) = -\sum_i p_i\log p_i$ subject to $\sum_i p_i = 1$.
Show the optimum is $p_i = 1/n$ for all $i$ and that the maximal value is $\log n$. (This re-derives the
upper bound of Theorem~\ref{ch33:entropy-bounds} by calculus rather than Jensen.)
\end{exercise}

\begin{exercise}[A Huffman-style code]
\label{ch33:ex-huffman}
A source emits symbols $A,B,C,D,E$ with probabilities $0.4, 0.2, 0.2, 0.1, 0.1$. (a) Compute the entropy.
(b) Build a prefix-free binary code by repeatedly merging the two least-probable symbols (Huffman's
algorithm) and report the average codeword length. (c) By how many bits per symbol does it beat the naive
$3$-bit fixed-length code, and how close is it to the entropy lower bound?
\end{exercise}

\begin{exercise}[KL is not symmetric]
\label{ch33:ex-kl-asym}
Let $p = (0.7, 0.3)$ and $q = (0.5, 0.5)$. Compute $\KL(p\|q)$ and $\KL(q\|p)$ in bits, confirm both are
positive, and confirm they differ. Then verify the identity $H(p,q) = H(p) + \KL(p\|q)$ numerically for
this pair.
\end{exercise}

\begin{exercise}[Cross-entropy as a loss]
\label{ch33:ex-crossent-loss}
A three-class classifier outputs probabilities $q = (0.7, 0.2, 0.1)$ on a particular example whose true
label is class $2$ (the middle one). (a) Write the one-hot true distribution $p$. (b) Compute the
cross-entropy loss $H(p,q)$ in bits. (c) Show that the loss equals $-\log_2 q_{\text{true}}$ and explain
why only the true-class probability appears. (d) What loss would a perfectly confident correct prediction
incur?
\end{exercise}

\begin{exercise}[Mutual information of a noisy channel]
\label{ch33:ex-bsc}
For the binary symmetric channel with flip probability $f$ and a fair input bit, show that $I(X;Y) = 1 -
H_b(f)$. Tabulate $I$ for $f \in \{0, 0.1, 0.25, 0.5\}$ and interpret the endpoints. At what $f$
is the channel useless, and why?
\end{exercise}

\begin{exercise}[Conditioning can increase specific uncertainty]
\label{ch33:ex-conditioning}
Construct an explicit joint distribution of two binary variables $X,Y$ for which there is a value $y$ with
$H(X\given Y=y) > H(X)$, yet the \emph{average} $H(X\given Y) \le H(X)$ still holds. Verify both
inequalities numerically. (This makes the paradox box concrete.)
\end{exercise}

\begin{exercise}[MLE equals minimum cross-entropy --- proof]
\label{ch33:ex-mle-proof}
Let $x_1,\dots,x_N$ be i.i.d.\ draws over a finite alphabet, with empirical distribution $\hat p$. Prove
that $\frac1N\sum_i \log q_\theta(x_i) = -H(\hat p, q_\theta)$, and conclude that maximizing the
log-likelihood over $\theta$ is equivalent to minimizing $\KL(\hat p\|q_\theta)$. State precisely which
term in the decomposition does not depend on $\theta$ and why that makes the two optimizations identical.
\end{exercise}

\begin{exercise}[Independence and mutual information]
\label{ch33:ex-mi-indep}
Let $X \in \{-1,0,1\}$ be uniform and let $Y = X^2$. (a) Compute $\Corr(X,Y)$ and observe it is zero.
(b) Compute $I(X;Y)$ and show it is strictly positive. (c) Explain in one sentence why mutual information
detects the dependence that correlation completely misses.
\end{exercise}

\begin{exercise}[Grouping / recursivity of entropy --- proof]
\label{ch33:ex-grouping}
Show that entropy satisfies the grouping property: if outcomes are partitioned into groups, the total
entropy equals the entropy of the group choice plus the average within-group entropy,
\[
  H(p_1,\dots,p_n) = H(\text{group probabilities}) + \sum_{\text{groups } g} \P(g)\,H(\text{conditional dist.\ within } g).
\]
Verify it on $(\tfrac12,\tfrac14,\tfrac14)$ grouped as $\{1\}$ and $\{2,3\}$, and explain how this is the
chain rule in disguise.
\end{exercise}
